\scanpage{42}
une \uline{application} \uline{composition}, et l'axiome à
formuler pour avoir une extension \deleted{en}
\inserted{des \(\Sigma_i,\Sigma_{ij}\)}
en des \(\Sigma_{i_0\ldots i_n}\), est que cette composition
(qui est compatible avec les unités, \unsure{manifestement},
par construction) soit une structure de
\uline{\(\mathcal M\)-catégorie}\(^*\). \deleted{\unsure{une}}
Donc la donnée de \(\Sigma\) satisfaisant (3)
équivaut à celle d'un Cat-objet de \(\mathcal M\),
au dessus de la catégorie \unsure{discrète} \(I_{\mathcal M}\) :
\[
 \Sigma_{\mathcal M}\longrightarrow I_{\mathcal M}.
\]

\margin{\(^*\) En fait, c'est un objet
\uline{ordonné} de \(\mathcal M\), et
\(\Sigma_{i_0\ldots i_n}\) s'identifie à l'ens des drapeaux
\[
 x_{i_0}\le x_{i_1}\le\cdots\le x_{i_n}
\]
de \(\Sigma_*\), au dessus des drapeaux
\(i_0-\cdots-i_n\) de \(I\).
Donc pour vérifier \unsure{(cf. p. 15, marge)}
\[
 \Sigma_{i_0\ldots i_n}
 =\Sigma_{\unsure{i_0i_1|i_0i_n}}\cap\cdots
 \cap\Sigma_{\unsure{i_0i_n}}
\]
il suffit de le vérifier pour \(n=2\), i.e.
\[
 \Sigma_{\unsure{ijk}}=\Sigma_{ik}\cap\Sigma_{jk},
\]
ce qui signifie que pour
tt \(z\in\Sigma_k\), \(\Sigma_*\unsure{-|\longrightarrow}I\) induit
\diagram{\begin{tikzcd}[column sep=large]
 \Sigma_*(z)\arrow[r,hook]\arrow[d,equal]&I(k)\arrow[d,equal]\\
 \{x\in\Sigma_*\mid x\le z\}&\{i\in I\mid i\le k\}
\end{tikzcd}}
qui est tel que l'ordre de \(\Sigma(z)\)
soit \uline{induit} par celui de \(I(k)\).
\deleted{Cela signifie}}

A l'aide de ces données, je vais
construire un espace stratifié de type \(I\),
sans utiliser en fait \deleted{aucune des hypothèses}
\inserted{les \deleted{\unsure{deux}} conditions}
(1) \deleted{et (2)}, la condition (2) n'intervenant
que sous la forme plus faible que
\[
 \Sigma_{i_0\ldots i_n}
 -\Sigma_{i_0\ldots\widehat{i_p}\ldots i_n}\qquad(p\ne n)
\]
est un plongement \deleted{bord}
(\unsure{si on ne veut pas construire} les structures côniques
relatives à l'espace stratifié obtenu).

Je considère pour ceci les couples
\deleted{\(d,d'\)} \((d\subset d')\) de deux drapeaux,
\deleted{\unsure{une}} tels que
\(d\) \deleted{de \(d'\)} ne soit pas cofinal dans \(d'\).
(Pour \(d'\) \inserted{\(=\{i_0,\ldots,i_n\}\)}
\deleted{de longueur \(n\)}
il y a donc \(2^n\) choix possibles pour \(d\)).

\margin{41}
